% ==================================================
% ABSTRACT
% ==================================================

\chapter*{Abstract}
\addcontentsline{toc}{chapter}{Abstract}

Hostel rooms are allotted once, and a student who gets the wrong hostel,
floor, direction, room type or roommate has no structured way to change
it. A direct one-to-one swap rarely works, because two students seldom
want exactly what the other has. A longer rotation usually does work
(A takes B's room, B takes C's, C takes A's), but no single student can
see such a chain, because seeing it needs every student's room and every
student's preferences at once.

The \textbf{Hostel Room Swapping Service} holds that information
centrally and computes the chains. Because the only rooms available are
the ones the participants already hold, a solution is a permutation of
students over bed slots, and the swap chains are that permutation's
cycles. The engine builds a students-by-slots cost matrix from weighted
hard and soft preferences, solves it with the Hungarian algorithm, uses
Murty's algorithm to enumerate the next best assignments, splits each
assignment into its cycles, and keeps only the chains in which nobody is
worse off and somebody is better off. Each student is shown up to three
such chains, and a chain moves rooms only after every member approves it
within three days.

The final system is a FastAPI backend over PostgreSQL (Supabase), with a
React and TypeScript interface for students and wardens, session-cookie
authentication, an append-only allocation history enforced by database
triggers, and a schema managed as numbered migrations. On the frozen
30-student evaluation cohort, one matching round raises mean preference
satisfaction from 11\% to 77\% with nobody worse off; searching 50
alternative assignments instead of one cuts the longest chain that must
be approved together from 28 students to 11 at no cost in satisfaction.
A benchmark over ten seeded cohorts of 300 students set the chain-length
cap at eight, which reaches 58\% of students with an executable offer.
Murty's search was later made 4--13 times faster, with identical
results, by solving each subproblem on a smaller matrix. The code base
is covered by 152 automated tests, four of them against PostgreSQL, run
in continuous integration on
Python 3.10--3.12.

\vspace{1em}
\noindent\textbf{Keywords:} assignment problem, Hungarian algorithm,
Murty's algorithm, exchange cycles, Pareto improvement, hostel
allocation, state machine, software engineering

% ==================================================
% ACKNOWLEDGEMENT
% ==================================================

\chapter*{Acknowledgement}
\addcontentsline{toc}{chapter}{Acknowledgement}

We thank our supervisor, Dr.\ Raghav B.\ Venkataramaiyer, for his
guidance and feedback throughout the semester, and in particular for the
questions about chain invalidation that shaped the approval workflow.
We also thank the Computer Science and Engineering Department of Thapar
Institute of Engineering and Technology for the opportunity to build a
complete software system as part of UCS503 Software Engineering.

We are grateful to the authors and maintainers of the open-source
projects this service is built on, in particular SciPy, FastAPI,
psycopg, PostgreSQL, React, Vite, pytest and Material for MkDocs.

\vspace{2cm}
\begin{flushright}
Aabhas Khandelwal (\texttt{1024030673}) \\
Aakash Chawla (\texttt{1024031097}) \\
Joy Soni (\texttt{1024030674})
\end{flushright}
